\documentclass[10pt,journal]{IEEEtran}
\usepackage{booktabs}
\usepackage{hyperref}
\usepackage{microtype}

\title{Agraya: A Secure Agentic Platform for Farmer Decision Support}
\author{Agri-Agentic Suite Contributors}

\begin{document}
\maketitle

\begin{abstract}
Agraya is an end-to-end agricultural platform that combines a farmer marketplace,
crop and field image analysis, grounded advisory services, farm planning, curated
news, and Telegram access. A Next.js application presents one coherent workflow,
while a FastAPI gateway protects and coordinates specialized local-AI services.
Supabase supplies Google authentication, PostgreSQL, object storage, and row-level
security. Ollama keeps language and vision inference deployable on farmer-focused
infrastructure without depending on a paid model API.
\end{abstract}

\section{System Architecture}
The browser communicates only with the platform API and Supabase. Supabase access
tokens authenticate private API operations; ownership is also enforced in
PostgreSQL through row-level security (RLS). The platform API validates requests,
stores domain records, handles uploads, ingests news, receives Telegram webhooks,
and forwards AI work over a private Docker network. Specialized services remain
independently replaceable: GroundTruth performs retrieval-augmented question
answering, YieldWise produces financial plans, and FieldScout runs multimodal crop
diagnosis through Ollama.

\begin{table}[h]
\caption{Primary Components}
\centering
\begin{tabular}{p{0.25\linewidth}p{0.67\linewidth}}
\toprule
Component & Responsibility \\
\midrule
Next.js UI & Public product pages and authenticated farmer workflows. \\
Platform API & Authentication, validation, persistence, orchestration, and integrations. \\
AI services & RAG advice, farm finance, vision diagnosis, and speech responses. \\
Supabase & Google Auth, PostgreSQL, Storage, RLS, and migration-managed schema. \\
Ollama & Local text and vision model runtime. \\
Telegram & Image intake and delivery of diagnosis results. \\
\bottomrule
\end{tabular}
\end{table}

Deployment is container-first. Docker Compose starts the web app, gateway, three
AI services, and Ollama. Persistent volumes retain models, embeddings, and trained
artifacts. Root-level pnpm commands and path-safe scripts provide the same entry
point on Linux, macOS, and Windows Docker Desktop.

\newpage
\section{Use Cases and Feature Set}
The principal actor is a farmer who signs in with Google, records one or more
farms, publishes produce, requests advice, and reviews saved diagnoses. Buyers may
browse active listings without authentication and contact sellers. Administrators
can refresh trusted news feeds, while a linked Telegram account gives the farmer a
low-friction mobile capture channel.

\begin{table}[h]
\caption{End-to-End User Flows}
\centering
\begin{tabular}{p{0.27\linewidth}p{0.64\linewidth}}
\toprule
Flow & System behavior \\
\midrule
Marketplace & Validate produce data, upload an image, publish a listing, then expose active inventory publicly. \\
Crop doctor & Accept JPEG, PNG, or WebP; select disease, crop, or field mode; call FieldScout; persist structured findings. \\
Advisor & Answer agricultural questions from local documents, accept voice, or create a farm financial plan. \\
News hub & Normalize configured RSS/Atom sources, de-duplicate articles, and serve a current public feed. \\
Telegram & Link a chat using a short-lived token, process a photo, store the diagnosis, and reply in the chat. \\
\bottomrule
\end{tabular}
\end{table}

FieldScout returns a condition, confidence, severity, observations, causes,
recommended actions, and safety notes rather than unstructured prose alone. The
same analysis modes are shared by web and Telegram entry points. Marketplace,
farm, diagnosis, and Telegram records are owner-scoped. News and active listings
are intentionally public. The user interface exposes dedicated routes for the
marketplace, crop doctor, advisor, farms, news, diagnoses, Telegram integration,
and login, with a long-form product home page connecting these capabilities.

The first production iteration deliberately favors modular service boundaries.
Additional agents can be introduced behind the platform API without exposing new
internal ports or changing the browser's trust model. Likely extensions include
weather alerts, localized crop calendars, multilingual output, buyer messaging,
price notifications, and human agronomist escalation.

\newpage
\section{Security, Reliability, and Operation}
Security is layered. Supabase validates Google identities, and the API verifies
their bearer tokens before private operations. RLS uses \texttt{auth.uid()} for
record ownership; user-editable metadata is not used for authorization. Public
clients receive only the Supabase publishable key. The Supabase secret key,
database URL, Telegram token, webhook secret, and internal API key remain
server-only. Internal AI endpoints require a shared header and are not intended
for browser use.

Uploads have bounded sizes and permitted MIME types. Filenames are generated by
the server, temporary audio files are cleaned up, external requests use timeouts,
and SQL access is parameterized. CORS is restricted to configured origins.
Container health checks and restart policies support recovery, while Ollama and
embedding volumes avoid repeated downloads. Telegram update identifiers are
stored to prevent duplicate processing. A production deployment should add a
durable job queue for webhook work, centralized telemetry, rate limiting, backups,
and periodic Supabase security and performance advisor reviews.

\section{Data and Development Lifecycle}
The migration defines profiles, farms, listings, diagnoses, news, Telegram links,
commodity prices, indexes, storage buckets, grants, and RLS policies. It should be
reviewed and applied from a controlled branch before production. Developers copy
\texttt{.env.example} to \texttt{.env}, install pinned JavaScript dependencies
with pnpm, and launch the complete system with \texttt{pnpm run dev}. Python
dependencies are locked for uv-based development and container builds.

\section{Conclusion}
Agraya turns previously separate agricultural prototypes into one secure product
surface. Its gateway-centered architecture, owner-scoped data model, local model
runtime, and web/Telegram parity provide a practical base for field pilots while
preserving clear paths toward scale and human-in-the-loop support.

\begin{thebibliography}{1}
\bibitem{supabase} Supabase, ``Platform Documentation,'' 2026. [Online]. Available: \url{https://supabase.com/docs}
\bibitem{ollama} Ollama, ``Ollama Documentation,'' 2026. [Online]. Available: \url{https://docs.ollama.com}
\end{thebibliography}
\end{document}
